\documentclass[11pt,a4paper]{report}
$if(fontfamily)$
\usepackage[$for(fontfamilyoptions)$$fontfamilyoptions$$sep$,$endfor$]{$fontfamily$}
$else$
\usepackage{lmodern}
$endif$
$if(linestretch)$
\usepackage{setspace}
\setstretch{$linestretch$}
$endif$
\usepackage{amssymb,amsmath}
\usepackage{ifxetex,ifluatex}
\usepackage{fixltx2e} % provides \textsubscript
\ifnum 0\ifxetex 1\fi\ifluatex 1\fi=0 % if pdftex
  \usepackage[$if(fontenc)$$fontenc$$else$T1$endif$]{fontenc}
  \usepackage[utf8]{inputenc}
$if(euro)$
  \usepackage{eurosym}
$endif$
\else % if luatex or xelatex
  \ifxetex
    \usepackage{mathspec}
  \else
    \usepackage{fontspec}
  \fi
  \defaultfontfeatures{Ligatures=TeX,Scale=MatchLowercase}
  \newcommand{\euro}{€}
$if(mainfont)$
    \setmainfont[$for(mainfontoptions)$$mainfontoptions$$sep$,$endfor$]{$mainfont$}
$endif$
$if(sansfont)$
    \setsansfont[$for(sansfontoptions)$$sansfontoptions$$sep$,$endfor$]{$sansfont$}
$endif$
$if(monofont)$
    \setmonofont[Mapping=tex-ansi$if(monofontoptions)$,$for(monofontoptions)$$monofontoptions$$sep$,$endfor$$endif$]{$monofont$}
$endif$
$if(mathfont)$
    \setmathfont(Digits,Latin,Greek)[$for(mathfontoptions)$$mathfontoptions$$sep$,$endfor$]{$mathfont$}
$endif$
$if(CJKmainfont)$
    \usepackage{xeCJK}
    \setCJKmainfont[$for(CJKoptions)$$CJKoptions$$sep$,$endfor$]{$CJKmainfont$}
$endif$
\fi
% use upquote if available, for straight quotes in verbatim environments
\IfFileExists{upquote.sty}{\usepackage{upquote}}{}
% use microtype if available
\IfFileExists{microtype.sty}{%
\usepackage{microtype}
\UseMicrotypeSet[protrusion]{basicmath} % disable protrusion for tt fonts
}{}
$if(geometry)$
\usepackage[$for(geometry)$$geometry$$sep$,$endfor$]{geometry}
$endif$
\usepackage{hyperref}
\PassOptionsToPackage{usenames,dvipsnames}{color} % color is loaded by hyperref
\hypersetup{unicode=true,
$if(title-meta)$
            pdftitle={$title-meta$},
$endif$
$if(author-meta)$
            pdfauthor={$author-meta$},
$endif$
$if(keywords)$
            pdfkeywords={$for(keywords)$$keywords$$sep$; $endfor$},
$endif$
$if(colorlinks)$
            colorlinks=true,
            linkcolor=$if(linkcolor)$$linkcolor$$else$Maroon$endif$,
            citecolor=$if(citecolor)$$citecolor$$else$Blue$endif$,
            urlcolor=$if(urlcolor)$$urlcolor$$else$Blue$endif$,
$else$
            pdfborder={0 0 0},
$endif$
            breaklinks=true}
\urlstyle{same}  % don't use monospace font for urls
$if(lang)$
\ifnum 0\ifxetex 1\fi\ifluatex 1\fi=0 % if pdftex
  \usepackage[shorthands=off,$for(babel-otherlangs)$$babel-otherlangs$,$endfor$main=$babel-lang$]{babel}
$if(babel-newcommands)$
  $babel-newcommands$
$endif$
\else
  \usepackage{polyglossia}
  \setmainlanguage[$polyglossia-lang.options$]{$polyglossia-lang.name$}
$for(polyglossia-otherlangs)$
  \setotherlanguage[$polyglossia-otherlangs.options$]{$polyglossia-otherlangs.name$}
$endfor$
\fi
$endif$
$if(natbib)$
\usepackage{natbib}
\bibliographystyle{$if(biblio-style)$$biblio-style$$else$plainnat$endif$}
$endif$
$if(biblatex)$
\usepackage{biblatex}
$for(bibliography)$
\addbibresource{$bibliography$}
$endfor$
$endif$
$if(listings)$
\usepackage{listings}
$endif$
$if(lhs)$
\lstnewenvironment{code}{\lstset{language=Haskell,basicstyle=\small\ttfamily}}{}
$endif$
$if(highlighting-macros)$
$highlighting-macros$
$endif$
$if(verbatim-in-note)$
\usepackage{fancyvrb}
\VerbatimFootnotes % allows verbatim text in footnotes
$endif$
$if(tables)$
\usepackage{longtable,booktabs}
$endif$
$if(graphics)$
\usepackage{graphicx,grffile}
\makeatletter
\def\maxwidth{\ifdim\Gin@nat@width>\linewidth\linewidth\else\Gin@nat@width\fi}
\def\maxheight{\ifdim\Gin@nat@height>\textheight\textheight\else\Gin@nat@height\fi}
\makeatother
% Scale images if necessary, so that they will not overflow the page
% margins by default, and it is still possible to overwrite the defaults
% using explicit options in \includegraphics[width, height, ...]{}
\setkeys{Gin}{width=\maxwidth,height=\maxheight,keepaspectratio}
$endif$
$if(links-as-notes)$
% Make links footnotes instead of hotlinks:
\renewcommand{\href}[2]{#2\footnote{\url{#1}}}
$endif$
$if(strikeout)$
\usepackage[normalem]{ulem}
% avoid problems with \sout in headers with hyperref:
\pdfstringdefDisableCommands{\renewcommand{\sout}{}}
$endif$
$if(indent)$
$else$
\setlength{\parindent}{0pt}
\setlength{\parskip}{6pt plus 2pt minus 1pt}
$endif$
\setlength{\emergencystretch}{3em}  % prevent overfull lines
\providecommand{\tightlist}{%
  \setlength{\itemsep}{0pt}\setlength{\parskip}{0pt}}
$if(numbersections)$
\setcounter{secnumdepth}{5}
$else$
\setcounter{secnumdepth}{0}
$endif$
$if(subparagraph)$
$else$
% Redefines (sub)paragraphs to behave more like sections
\ifx\paragraph\undefined\else
\let\oldparagraph\paragraph
\renewcommand{\paragraph}[1]{\oldparagraph{#1}\mbox{}}
\fi
\ifx\subparagraph\undefined\else
\let\oldsubparagraph\subparagraph
\renewcommand{\subparagraph}[1]{\oldsubparagraph{#1}\mbox{}}
\fi
$endif$
$if(dir)$
\ifxetex
  % load bidi as late as possible as it modifies e.g. graphicx
  $if(latex-dir-rtl)$
  \usepackage[RTLdocument]{bidi}
  $else$
  \usepackage{bidi}
  $endif$
\fi
\ifnum 0\ifxetex 1\fi\ifluatex 1\fi=0 % if pdftex
  \TeXXeTstate=1
  \newcommand{\RL}[1]{\beginR #1\endR}
  \newcommand{\LR}[1]{\beginL #1\endL}
  \newenvironment{RTL}{\beginR}{\endR}
  \newenvironment{LTR}{\beginL}{\endL}
\fi
$endif$
$for(header-includes)$
$header-includes$
$endfor$

$if(title)$
\title{$title$}
$endif$
$if(subtitle)$
\providecommand{\subtitle}[1]{}
\subtitle{$subtitle$}
$endif$
$if(author)$
\author{$for(author)$$author$$sep$ \and $endfor$}
$endif$

%%%%%%%%%%%%%%%%%%%%%%%%%%%%%%%%%%%%%%%%%
% ADD YOUR \usepackage{...} COMMANDS HERE
%%%%%%%%%%%%%%%%%%%%%%%%%%%%%%%%%%%%%%%%%


\newcommand{\thema}{[Thema der Bachelorarbeit]}
\newcommand{\schlagworte}{[Platz, f\"ur, spezifische, Schlagworte, zur, Ausarbeitung ]}
\newcommand{\zusammenfassung}{[Text der Zusammenfassung etwa 150 Worte. Es soll der
  L"osungsweg beschrieben sein.]}
\newcommand{\ausgabedatum}{[Datum]}
\newcommand{\abgabedatum}{[Datum]}
\newcommand{\autor}{[Vor- und Zuname des/der Bachelorkandidaten/in]}
\newcommand{\autorStrasse}{[Stra"se]}
\newcommand{\autorPLZ}{[PLZ]}
\newcommand{\autorOrt}{[Ort]}
\newcommand{\autorGeburtsort}{[Geburtsort]}
\newcommand{\autorGeburtsdatum}{[Datum]}
\newcommand{\prueferA}{[Titel, Vor- und Zuname des 1. Pr"ufers]}
\newcommand{\prueferB}{[Titel, Vor- und Zuname des 2. Pr"ufers]}
\newcommand{\firma}{[HTWG oder Firmenname]}
\newcommand{\studiengang}{[Software-Engineering/Technische Informatik/Wirtschaftsinformatik]}

\date{$date$}

\begin{document}


\begin{titlepage}

\vspace*{-3.5cm}

\begin{flushleft}
\hspace*{-1cm} \includegraphics[width=15.7cm]{htwg-logo}
\end{flushleft}

\vspace{2.5cm}

\begin{center}
	\huge{
		\textbf{\thema} \\[5cm]
	}
	\Large{
		\textbf{\autor}} \\[6.5cm]
	\large{
		\textbf{Konstanz, \abgabedatum} \\[2.3cm]
	}
	
	\Huge{
		\textbf{{\sf BACHELORARBEIT}}
	}
\end{center}

\end{titlepage}

\thispagestyle{empty}
{
\setlength{\parskip}{0.5cm}
        \begin{center}
        \textbf{\huge BACHELORARBEIT}

        \textbf{zur Erlangung des akademischen Grades}

        \textbf{\Large Bachelor of Science (B. Sc.)}

        \textbf{an der}

        \textsf{\huge Hochschule Konstanz}\\
        {\small Technik, Wirtschaft und Gestaltung}

        \textsf{\Large Fakult"at Informatik} \\
        Studiengang \studiengang
        \end{center}
}
\begin{center}

\vspace*{2cm}

\begin{tabular}{p{3cm}p{10cm}}
Thema: & \multicolumn{1}{l}{\textbf{\large \thema}} \\[15ex]
Bachelorkandidat: & \autor, \autorStrasse, \autorPLZ{}  \autorOrt{} \\[15ex]
1. Pr"ufer: & \prueferA \\
2. Pr"ufer: & \prueferB \\[25ex]
Ausgabedatum: & \ausgabedatum \\
Abgabedatum: & \abgabedatum \\
\end{tabular}
\end{center}

\begin{center}
{\Large \textbf{Zusammenfassung (Abstract)}}
\end{center}

\bigskip

\begin{center}
	\begin{tabular}{p{2.8cm}p{10cm}}
		Thema: & \thema \\
		 & \\
		Bachelorkandidat: & \autor \\
		 & \\
		Firma: & \firma \\
		 & \\
		Betreuer: & \prueferA  \\[.5ex]
		 &  \prueferB \\
		 & \\
		Abgabedatum: & \abgabedatum \\
		 & \\
		Schlagworte: & \schlagworte \\
		 & \\
	\end{tabular}
\end{center}

\bigskip

\noindent
\zusammenfassung
\chapter*{Ehrenw"ortliche Erkl"arung}
\addcontentsline{toc}{chapter}{Ehrenw"ortliche Erkl"arung}

Hiermit erkl"are ich
\textit{\autor, geboren am \autorGeburtsdatum{} in \autorGeburtsort{}}, dass ich\\

\begin{tabular}{lp{12cm}}
(1) & meine Bachelorarbeit mit dem Titel \\[1em]
& \textbf{\thema} \\[1em]
& selbstst"andig und ohne fremde Hilfe angefertigt und keine anderen als die angef"uhrten Hilfen benutzt habe;\\[1em]
(2) & die "Ubernahme w"ortlicher Zitate, von Tabellen, Zeichnungen, Bildern und
Programmen aus der Literatur oder anderen Quellen (Internet) sowie die Verwendung
der Gedanken anderer Autoren an den entsprechenden Stellen innerhalb der Arbeit
gekennzeichnet habe.\\
\end{tabular}

\vspace*{1cm}

\noindent
Ich bin mir bewusst, dass eine falsche Erkl"arung rechtliche Folgen haben wird.\\

\vspace*{3cm}

\noindent
Konstanz, \abgabedatum \hfill \begin{tabular}{c} \\ \\ \rule{5cm}{1pt} \\ (Unterschrift)\end{tabular}


\tableofcontents

$if(title)$
\maketitle
$endif$
$if(abstract)$
\begin{abstract}
$abstract$
\end{abstract}
$endif$

$for(include-before)$
$include-before$

$endfor$
$if(toc)$
{
$if(colorlinks)$
\hypersetup{linkcolor=$if(toccolor)$$toccolor$$else$black$endif$}
$endif$
\setcounter{tocdepth}{$toc-depth$}
\tableofcontents
}
$endif$
$if(lot)$
\listoftables
$endif$
$if(lof)$
\listoffigures
$endif$
$body$

$if(natbib)$
$if(bibliography)$
$if(biblio-title)$
$if(book-class)$
\renewcommand\bibname{$biblio-title$}
$else$
\renewcommand\refname{$biblio-title$}
$endif$
$endif$
\bibliography{$for(bibliography)$$bibliography$$sep$,$endfor$}

$endif$
$endif$
$if(biblatex)$
\printbibliography$if(biblio-title)$[title=$biblio-title$]$endif$

$endif$
$for(include-after)$
$include-after$

$endfor$
\end{document}
